% =========================================================================
\chapter{Edge Computing \& Specialized FPGA Hardware Accelerators}
\label{ch:edge_specialized}
% =========================================================================

\section{Status Quo \& Key Challenges in Edge Acceleration}
As AI accelerators transition from datacenters to real-time edge environments (smart mobility, industrial automation, drone navigation, edge vision), hardware engineers face the fundamental challenge of balancing **compute performance, power consumption, silicon die cost, and thermal limits**.

In a comprehensive overview of FPGA-based CNN accelerators for edge computing, Fata et al. (IEEE Access 2025)~\cite{fata2025balancing} noted that the edge AI market exceeded $\$6\,\text{Billion}$ with a $30.4\%$ CAGR. However, edge deployment is severely constrained by three primary limitations:
\begin{enumerate}[leftmargin=*]
    \item \textbf{Memory Bandwidth Bottleneck}: Fetching large neural network weight parameters from off-chip DRAM stalls execution datapaths.
    \item \textbf{Heterogeneous Model Complexity}: Edge applications combine diverse network layers (Convolutions, Recurrent LSTM/GRU units, Transformers, and custom post-processing).
    \item \textbf{Strict Die Area \& Cost Bounds}: Edge hardware must utilize low-cost FPGAs or compact ASIC dies with limited Logic Lookup Tables (LUTs), DSP Slices, and Block RAM (BRAM).
\end{enumerate}

To address these challenges, researchers have developed specialized silicon hardware frameworks and algorithmic optimizations.

% -------------------------------------------------------------------------
\section{Reconfigurable CNN-RNN Frameworks on FPGA}
% -------------------------------------------------------------------------
Many real-time applications (such as speech recognition, autonomous driving trajectory forecasting, and video action analysis) combine spatial Convolutional Neural Networks (CNNs) with temporal Recurrent Neural Networks (RNNs/LSTMs).

Zeng et al. (Tsinghua University / DeePhi Technology)~\cite{zeng2018efficient} designed a complete reconfigurable hardware framework for deploying general-purpose hybrid CNN-RNN models on FPGAs.

Key innovations of the DeePhi framework:
\begin{itemize}[leftmargin=*]
    \item \textbf{Aristotle \& Descartes IP Cores}: Integrated DeePhi's Aristotle IP (optimized for 2D CONV spatial feature extraction) and Descartes IP (optimized for sparse matrix-vector RNN temporal processing) into a single reconfigurable hardware system.
    \item \textbf{Co-Optimization Toolchain}: Developed a hardware-software compiler co-optimization toolchain integrated with the Xilinx SDSoC environment. The toolchain automatically identifies optimal IP tile configurations, maximizing FPGA resource utilization.
    \item \textbf{Quantization and Pruning}: Applied weight pruning and 8-bit fixed-point quantization, achieving over **$4\times$ reduction in model size** with zero loss in recognition accuracy.
\end{itemize}

% -------------------------------------------------------------------------
\section{DDR-Less Stereo Matching for Co-Axial RGB-NIR Sensors}
% -------------------------------------------------------------------------
Edge vision robots require all-day, 24/7 depth perception across changing lighting conditions (bright daylight, dusk, and complete darkness). Co-axial RGB-Near-Infrared (RGB-NIR) sensor fusion provides robust multi-modal perception.

However, traditional stereo matching algorithms rely heavily on external DDR memory buffering to perform dynamic vertical residual compensation, stalling continuous hardware dataflow and violating microsecond latency bounds.

To solve this, Zhang et al. (ACM GLSVLSI 2026)~\cite{zhang2026fpga} presented an **adaptive texture-aware DDR-less stereo matching accelerator** tailored for coaxial RGB-NIR sensors.

\begin{hardwarebox}[DDR-Less Hardware Datapath Architecture]
Zhang et al. eliminated off-chip DDR memory entirely by implementing a fully pipelined datapath with on-chip line buffers:
\begin{enumerate}[leftmargin=*]
    \item **On-Chip Line Buffer Grid**: Stores local image scanlines inside high-speed BRAM, enabling continuous stream processing as pixels arrive from the image sensor.
    \item **Adaptive Texture-Aware Cost Aggregation**: Computes census transforms and cross-modal correlation coefficients in parallel across RGB and NIR channels.
    \item **Zero DDR Memory Stalls**: Bypassing off-chip DDR transfers reduced depth estimation latency to **sub-millisecond microsecond bounds** while reducing hardware power consumption by **$62\%$**.
\end{enumerate}
\end{hardwarebox}

% -------------------------------------------------------------------------
\section{Hardware Accelerators for Homomorphic Encryption \& 5G NR}
% -------------------------------------------------------------------------
Beyond deep learning, edge physical AI hardware integrates specialized cryptographic and wireless telecommunication processing engines.

\subsection{Leveled Ring-LWE Fully Homomorphic Encryption (FHE)}
Privacy-preserving Physical AI requires executing mathematical inference directly over encrypted data without decrypting it. 

Su et al. (IEEE Access 2020)~\cite{su2020fpga} designed an FPGA hardware accelerator for **Leveled Ring-Learning With Errors (Ring-LWE) Fully Homomorphic Encryption**. 

Su et al. developed high-speed hardware polynomial multiplication engines utilizing **Number Theoretic Transform (NTT)** pipelines. The FPGA accelerator achieved a **$31\times$ speedup** in homomorphic polynomial multiplication over CPU software solvers, enabling secure edge AI inference.

\subsection{Scalable 5G NR Rate-Matching Engines}
5G New Radio (NR) edge base stations require ultra-low latency channel coding. 

Filipović et al. (IEEE Access 2025)~\cite{filipovic2025scalable} presented a scalable **5G NR Rate-Matcher and Rate-Dematcher FPGA accelerator**. By implementing parallel circular buffer management and bit-selection logic, the accelerator achieved 5G NR peak data rate requirements ($>10\,\text{Gbit/s}$) with minimal FPGA logic utilization.

% -------------------------------------------------------------------------
\section{Automatic Source Code Annotation \& Compiler Hints}
% -------------------------------------------------------------------------
Deploying complex mathematical algorithms (such as Newton's method or matrix factorizations) onto high-performance CPU, GPU, or FPGA hardware requires inserting compiler hints and pragmas (such as \texttt{\#pragma HLS UNROLL}, \texttt{\#pragma HLS PIPELINE}, OpenMP annotations) into C/C++ source code.

Inserting annotations manually requires extensive hardware design expertise and lacks code portability across different hardware targets.

To solve this, Shahzad et al. (Boston University / Red Hat / IEEE Access 2025)~\cite{shahzad2025annotationgym} developed **AnnotationGym**, a generic framework for automatic source code annotation.

Key capabilities of AnnotationGym:
\begin{itemize}[leftmargin=*]
    \item \textbf{Automated Exploration}: Uses machine learning search agents to explore optimal combinations of compiler pragmas and loop unrolling factors.
    \item \textbf{Hardware Performance Tuning}: Automatically optimizes C/C++ code for Xilinx High-Level Synthesis (HLS) and Intel FPGA compilers, achieving up to **$3.2\times$ performance improvement** over human-annotated baseline code.
\end{itemize}
